% !PRINTED 299
\weberpage{299}
%<<CONT
into $[z_1+\alpha_1, z_2+\alpha_2, \ldots, z_k+\alpha_k]$. Accordingly we can also denote the permutation $A$ thus:
\weq{4}{\begin{pmatrix} z_1, & z_2, & \ldots, & z_k \\ z_1+\alpha_1, & z_2+\alpha_2, & \ldots, & z_k+\alpha_k \end{pmatrix},}
and one obtains the whole group $Q$, if one lets $\alpha_1, \alpha_2, \ldots, \alpha_k$ each run through a complete system of residues to the modulus $p$.

The group $Q$, whose mode of formation is now established, must be a normal divisor of $P$, and from this a general form can be derived in which all permutations of $P$ must be contained.

For this we require a lemma on the analytic representation of the permutations, which we first derive.

\begin{center}
\S.\,77.\\
\gesperrt{Analytic Representation of the Permutations.}
\end{center}

\gesperrt{Lemma. If under any permutation of the quantities (2), \S.\,76, $z_1$ passes into $z_1'$, $z_2$ into $z_2'$, \ldots, $z_k$ into $z_k'$, then one can always set
\weq{1}{
\begin{aligned}
z_1' &\equiv \varphi_1(z_1, z_2, \ldots, z_k)\\
z_2' &\equiv \varphi_2(z_1, z_2, \ldots, z_k)\\
&\;\cdots\cdots\cdots\cdots\\
z_k' &\equiv \varphi_k(z_1, z_2, \ldots; z_k)
\end{aligned}
\quad (\mathrm{mod}\ p).
}
where $\varphi_1, \varphi_2, \ldots, \varphi_k$ denote integral rational functions of the variables $z_1, z_2, \ldots, z_k$ with integral coefficients, which do not exceed the degree $p - 1$ with respect to any of the variables.}

The theorem is a generalisation of the corresponding theorem in \S.\,180 of the first volume. There the theorem is proved for $k = 1$. We preface the general proof, which rests on complete induction, with the following considerations.

Since in (1) we have to put for the arguments $z_i$ all combinations of $k$ sequences to the modulus $p$, there result $p^k$ such systems. Each of the functions $\varphi$ has $p^k$ undetermined coefficients, and if therefore the $z_k'$ are given, we obtain $kp^k$ linear congruences for just as many unknown numbers, namely the coefficients of the $\varphi_k$. It remains only to show
%CONT>>
